\begin{exbox}[依次通过两个腔：总损失相同，光学记录可以不同]
保留两个互相独立、同频的模式 $a_1,a_2$，输入都是真空。
电子沿用同一个双边格，两个脉冲的交换参数为 $g_1,g_2$。
假定脉冲间的反冲色散相位可忽略，令
$X_j=g_jb^\dagger a_j-g_j^*ba_j^\dagger$。
求电子总损失的分布，以及给定总损失后两个腔的联合光态；取 $\mu=|g_1|^2+|g_2|^2>0$，使非真空结果有可能发生。

\textbf{解。} 不同模式算符对易，且 $bb^\dagger=b^\dagger b=I$，故 $[X_1,X_2]=0$。
两个相互作用可以先分别正规排序，再作用于真空，得到
\begin{equation}
|\Psi\rangle=\ee^{-\mu/2}\sum_{m_1,m_2\ge0}
\frac{\zeta_1^{m_1}\zeta_2^{m_2}}{\sqrt{m_1!m_2!}}
|-(m_1+m_2),m_1,m_2\rangle,
\quad\zeta_j=-g_j^*,\quad\mu=|g_1|^2+|g_2|^2.
\label{eq:feqo-two-cavity-vacuum-joint}
\end{equation}
两个光子数的无条件联合计数概率是两个独立 Poisson 分布；
电子却只记录总损失 $n=m_1+m_2$，其概率为 $\ee^{-\mu}\mu^n/n!$。
\begin{equation}
|\psi_{12|n}\rangle=\sum_{m=0}^n
\sqrt{\binom nm}\left(\frac{\zeta_1}{\sqrt\mu}\right)^m
\left(\frac{\zeta_2}{\sqrt\mu}\right)^{n-m}|m,n-m\rangle.
\label{eq:feqo-two-cavity-heralded-binomial}
\end{equation}
这是总光子数为 $n$ 的纯态，数分布为二项分布。
归一化由二项定理给出；相对复相位来自两个腔的耦合相位。
$n=1$ 且两腔耦合模长相同时，电子后选择制备
$(\ee^{\ii\arg\zeta_1}|1,0\rangle+
\ee^{\ii\arg\zeta_2}|0,1\rangle)/\sqrt2$，
在两个腔之间有一个量子比特的纠缠，尽管电子没有分辨光子来自哪一个腔。
若两腔间色散不可忽略，需在两个交换算符中间插入传播子；
不同交换路径将获得不同相位，不能继续使用上述可交换指数。
\end{exbox}

\begin{proposition}[条件光子计数与二阶关联]
设 $0<\eta=|g_1|^2/\mu<1$。
在上式条件态中，$\langle N_1\rangle=n\eta$、
$\langle N_2\rangle=n(1-\eta)$，且
\begin{equation}
\langle N_1N_2\rangle=n(n-1)\eta(1-\eta),\qquad
g_{12}^{(2)}=\frac{\langle N_1N_2\rangle}
{\langle N_1\rangle\langle N_2\rangle}=1-\frac1n\quad(n\ge1).
\label{eq:feqo-heralded-hbt-correlation}
\end{equation}
若不对电子后选择，两个模式的计数独立，$g_{12}^{(2)}=1$。
\end{proposition}
\begin{proof}
对二项变量 $M$，$\mathbb EM=n\eta$、$\mathbb E[M(M-1)]=n(n-1)\eta^2$。
由于 $N_2=n-M$，
$\mathbb E[M(n-M)]=n\mathbb EM-\mathbb EM^2
=n(n-1)\eta(1-\eta)$，得到归一化关联。
当 $n=0$ 时两个分母为零，归一化 $g^{(2)}$ 没有定义，不能写成零。
对单模数态 $|n\rangle$，同样由 $(a^\dagger)^2a^2=N(N-1)$ 得
$g^{(2)}(0)=1-1/n$（$n\ge1$）。
\end{proof}

\begin{exbox}[光子数尾部怎样控制电子计数核]
对一般数分布 $p_n$，已建立的前向核满足
$P_\ell=\sum_np_n\mathsf K_{\ell n}$、$0\le\mathsf K_{\ell n}\le1$。
求截断 $n>N$ 后的计数误差，并比较数态对角输入与相干输入在联合态层面的误差。

\textbf{解。} 保留 $n\le N$ 而不重新归一化，遗漏概率为 $\eta_N=\sum_{n>N}p_n$，于是
$0\le P_\ell-P_\ell^{(N)}\le\eta_N$。
若输入是数态对角混合，对截断输入重新归一化后，原输入与该输入的迹距离就是 $\eta_N$；
CPTP 演化和测量不会增大这个距离，因此任意完整计数分布的总变差也至多为 $\eta_N$。
\[
\eta_N^{\rm th}=\left(\frac{\bar n}{1+\bar n}\right)^{N+1},\qquad
\eta_N^{\rm coh}=1-\ee^{-\bar n}\sum_{n=0}^{N}\frac{\bar n^n}{n!}.
\]
例如平均数为一的热态，要求 $\eta_N<10^{-6}$ 时取 $N=19$ 即可。
对于相干输入，保留数态振幅并重新归一化得到纯态 $|\alpha_N\rangle$，
其重叠满足 $|\langle\alpha|\alpha_N\rangle|^2=1-\eta_N$，所以纯态迹距离是 $\sqrt{\eta_N}$。
单能电子的能谱只依赖输入数分布，仍可使用 $\eta_N$ 的人口核误差；
相位敏感的联合测量或条件态一般只能先使用 $\sqrt{\eta_N}$ 的态误差界。
这个界来自输入尾概率，不需要把未计算的每条高能响应核另作经验估计。
\end{exbox}
